\begin{frame}{Relações Comerciais Livres ou com Baixa Tarifa}
  % \begin{figure}
  %   \centering
  %   \includegraphics[width=0.5\textwidth]{imagens/grafo_denso.png}
  %   \caption{Exemplo de grafo denso}
  % \end{figure}
  \begin{itemize}
    \item \textbf{Análise de Grafos:} Grafos densos representam mercados com alta conectividade e poucas barreiras tarifárias.
    \item \textbf{Oportunidades Econômicas:} Facilita a formação de clusters de parceiros, promovendo integração regional e aumento do fluxo comercial.
    \item \textbf{Parcerias Não Desenvolvidas:} Identificação de vértices isolados ou pouco conectados pode indicar potenciais mercados a serem explorados.
    \item \textbf{Riscos:} Alta conectividade pode gerar dependência excessiva de determinados mercados ou vulnerabilidade a choques externos.
  \end{itemize}
\end{frame}

\begin{frame}{Relações Econômicas em Déficit e com Tarifa}
  % \begin{figure}
  %   \centering
  %   \includegraphics[width=0.5\textwidth]{imagens/grafo_esparso.png}
  %   \caption{Exemplo de grafo esparso}
  % \end{figure}
  \begin{itemize}
    \item \textbf{Análise de Grafos:} Grafos esparsos ou com arestas ponderadas negativamente representam barreiras tarifárias e baixo fluxo comercial.
    \item \textbf{Oportunidades:} Redução de tarifas pode transformar componentes desconexos em clusters colaborativos, estimulando novos negócios.
    \item \textbf{Clusters de Relações:} Identificação de subgrafos pode revelar grupos de países/setores que operam em déficit, sugerindo estratégias de integração.
    \item \textbf{Riscos:} Barreiras elevadas podem isolar mercados, dificultando parcerias e aumentando custos operacionais.
  \end{itemize}
\end{frame}

\begin{frame}{Relações Econômicas sem Tarifa}
  % \begin{figure}
  %   \centering
  %   \includegraphics[width=0.5\textwidth]{imagens/grafo_completo.png}
  %   \caption{Exemplo de grafo completo}
  % \end{figure}
  \begin{itemize}
    \item \textbf{Análise de Grafos:} Grafos completos ou quase completos indicam livre circulação de bens e serviços.
    \item \textbf{Oportunidades:} Maximização de rotas comerciais, diversificação de parceiros e aumento da resiliência econômica.
    \item \textbf{Observações:} Possibilidade de identificar ciclos de Hamilton para otimizar cadeias logísticas e rotas de distribuição.
    \item \textbf{Riscos:} Necessidade de monitorar clusters para evitar concentração excessiva e garantir equilíbrio competitivo.
  \end{itemize}
\end{frame}
